%=========================================================
%               CLASSE DU DOCUMENT
%=========================================================
\documentclass[12pt,a4paper]{article}

%=========================================================
%         ENCODAGE ET POLICES DE CARACTÈRES
%=========================================================
\usepackage[utf8]{inputenc}    % Encodage UTF-8 (inutile avec XeLaTeX/LuaLaTeX)
\usepackage[T1]{fontenc}       % Meilleure gestion des caractères accentués
\usepackage[french]{babel} % <--- Ce package fournit \og et \fg{}
\usepackage{lmodern}           % Police Latin Modern
\usepackage{times}             % Police Times (à supprimer si tu préfères Latin Modern)

%=========================================================
%                MATHÉMATIQUES
%=========================================================
\usepackage{amsmath}           % Environnements mathématiques avancés
\usepackage{amssymb}           % Symboles mathématiques supplémentaires
\usepackage{amsfonts}          % Polices mathématiques
\usepackage{mathrsfs}          % Police calligraphique (\mathscr)
\usepackage{tkz-tab}           % Tableaux de signes et de variations
\everymath{\displaystyle}      % Affiche toutes les formules en style "display"

%=========================================================
%               MISE EN PAGE
%=========================================================
\usepackage[
    left=1cm,
    right=1cm,
    top=1cm,
    bottom=1.7cm
]{geometry}                    % Marges du document

\usepackage{multicol}          % Texte sur plusieurs colonnes
\usepackage{float}             % Placement précis des figures et tableaux

%=========================================================
%            TABLEAUX ET LISTES
%=========================================================
\usepackage{array}             % Colonnes de tableaux personnalisées
\usepackage{multirow}          % Fusion de lignes dans les tableaux
\usepackage{enumitem}          % Personnalisation des listes

\usepackage{tikz}
\usetikzlibrary{angles,quotes,arrows.meta}
%=========================================================
%            COULEURS ET BOÎTES
%=========================================================
\usepackage[most]{tcolorbox}   % Boîtes colorées très personnalisables
\tcbuselibrary{skins,hooks}    % Bibliothèques supplémentaires de tcolorbox

\usepackage{mdframed}          % Cadres autour de paragraphes
\usepackage{varwidth}          % Largeur automatique des boîtes

%=========================================================
%                DESSINS ET GRAPHIQUES
%=========================================================
\usepackage{tikz}              % Dessins vectoriels
\usetikzlibrary{calc}
\usetikzlibrary{
    arrows,
    shapes.geometric,
    fit,
    patterns
}

\usepackage{pgfplots}          % Tracés de fonctions et graphiques
\pgfplotsset{compat=1.15}

%=========================================================
%              LIENS HYPERTEXTES
%=========================================================
\usepackage[
    colorlinks=true,
    linkcolor=blue,
    urlcolor=blue,
    citecolor=blue
]{hyperref}

%=========================================================
%          ÉLÉMENTS EN ARRIÈRE-PLAN
%=========================================================
\usepackage{eso-pic}           % Ajout de filigranes ou d'éléments sur chaque page

%=========================================================
%              SYMBOLES DIVERS
%=========================================================
\usepackage{pifont}            % Symboles (✓, ✗, etc.)

\DeclareUnicodeCharacter{2460}{\text{\ding{172}}} % Pour ①
\DeclareUnicodeCharacter{2461}{\text{\ding{173}}} % Pour ②
%=========================================================
%      PERSONNALISATION DES LISTES NUMÉROTÉES
%=========================================================

\newcommand\circitem[1]{%
\tikz[baseline=(char.base)]{
\node[circle,draw=gray,fill=red!55,
minimum size=1.2em,inner sep=0] (char) {#1};}}

\newcommand\boxitem[1]{%
\tikz[baseline=(char.base)]{
\node[fill=cyan,
minimum size=1.2em,inner sep=0] (char) {#1};}}

\setlist[enumerate,1]{label=\protect\circitem{\arabic*}}
\setlist[enumerate,2]{label=\protect\boxitem{\alph*}}

%=========================================================
%            COMMANDES PERSONNALISÉES
%=========================================================

% Soulignement coloré
\newcommand{\myul}[2][black]{%
\setulcolor{#1}\ul{#2}\setulcolor{black}}

% Tabulation horizontale
\newcommand\tab[1][1cm]{\hspace*{#1}}

%=========================================================
%             COMPTEURS PERSONNALISÉS
%=========================================================

%---------- Exemple ----------
\newcounter{exemple}

\newcommand{\exemple}{%
\refstepcounter{exemple}
\textbf{\textcolor{orange}{Exemple \theexemple : }}
\ignorespaces}

%---------- Solution ----------
\newcounter{solution}

\newcommand{\solution}{%
\refstepcounter{solution}
\textbf{\textcolor{orange}{Solution \thesolution : }}
\ignorespaces}

%---------- Exercices ----------
\definecolor{myorange}{rgb}{1.0,0.8,0.0}

\newcounter{exerciceapp}

\newcommand{\exerciceapp}{%
\refstepcounter{exerciceapp}
\textbf{\textcolor{myorange}
{Exercice d'application \theexerciceapp : }}
\ignorespaces}

%---------- Corrections ----------
\newcounter{correction}

\newcommand{\correction}{%
\refstepcounter{correction}
\textbf{\textcolor{myorange}
{Correction \thecorrection : }}
\ignorespaces}

%---------- Remarques ----------
\definecolor{myorange1}{rgb}{1.0,1.0,0.0}
\usepackage{enumitem}

\definecolor{myred}{RGB}{180, 40, 40}
\definecolor{myblue}{RGB}{20, 50, 120}
\newcounter{remarque}

\newcommand{\remarque}{%
\refstepcounter{remarque}
\textbf{\textcolor{myorange1}
{Remarque \theremarque : }}
\ignorespaces}

%=========================================================
%                 FILIGRANE
%=========================================================

\AddToShipoutPicture{
    \AtTextCenter{
        \makebox[0pt]{
            \rotatebox{45}{
                \textcolor[gray]{0.6}{
                    \fontsize{5cm}{5cm}\selectfont PGB
                }
            }
        }
    }
}

\begin{document}
% En-tête personnalisée
\begin{center}
    \Large\textbf{\underline{\textcolor{red}{Équations et Inéquations Du second dégré}}}\\[-0.1cm]
    \normalsize\textbf{Prof : M. BA} \hfill \textbf{Classe : 1erS2}\\[-0.1cm]
    \textbf{Année scolaire : 2025 -- 2026}
\end{center}
%=========================================================
% ========================================================================
\vspace{1em}
% PARTIE I : Trinome DU SECOND DEGRÉ
% ========================================================================

\section*{Objectifs spécifiques}
\begin{itemize}[label=$\checkmark$]
    \item Calculer \( (g \circ f)(x) \).
    \item Calculer \( (f \circ g)(x) \).
    \item Comprendre la non-commutativité de la composition.
\end{itemize}

\section*{Prérequis}
\begin{itemize}[label=$\checkmark$]
    \item Calcul dans \(\mathbb{R}\) (fonctions affines, carrées, polynômes).
    \item Notion de fonction (image, antécédent, ensemble de définition).
\end{itemize}

\section*{Supports didactiques}
\begin{itemize}[label=$\checkmark$]
    \item Collection N. Dimathème Terminale \( A_2 / A_3 \).
    \item Ordinateur (visualisations avec GeoGebra ou tableur).
\end{itemize}

\section*{Plan du chapitre}
\begin{enumerate}
    \item \textbf{Calcul de \( (g \circ f)(x) \)}
    \begin{enumerate}
        \item Exemples
        \begin{enumerate}
            \item Exemple 1
            \item Exemple 2
        \end{enumerate}
        \item Exercice d’application
    \end{enumerate}
    \item \textbf{Calcul de \( (f \circ g)(x) \)}
    \begin{enumerate}
        \item Exemples
        \begin{enumerate}
            \item Exemple 1
            \item Exemple 2
        \end{enumerate}
        \item Exercice d’application
    \end{enumerate}
\end{enumerate}

\newpage

\section*{Introduction}
La \textbf{composition d'applications} consiste à appliquer successivement deux fonctions. 
Si \( f \) transforme \( x \) en \( f(x) \) et \( g \) transforme \( y \) en \( g(y) \), alors \( g \circ f \) signifie \og on applique d'abord \( f \), puis \( g \) \fg{}.

\medskip
\noindent
\textbf{Notation :} 
\[
(g \circ f)(x) = g\big(f(x)\big)
\]

\section*{\color{myred} I. Calcul de \( (g \circ f)(x) \) }

\subsection*{\color{myred} 1. Exemples }
\subsubsection*{\color{myred} a. Exemple 1 (fonctions affines) }
Soit 
\[
f(x) = 2x + 1 \quad \text{et} \quad g(x) = 3x - 4.
\]
Calculons \( (g \circ f)(x) \) :
\[
(g \circ f)(x) = g(f(x)) = g(2x+1) = 3(2x+1) - 4 = 6x + 3 - 4 = 6x - 1.
\]

\subsubsection*{\color{myred} b. Exemple 2 (fonction carrée et racine carrée)}
Soit 
\[
f(x) = x^2 \quad \text{et} \quad g(x) = \sqrt{x} \quad (x \ge 0).
\]
Alors :
\[
(g \circ f)(x) = g(f(x)) = g(x^2) = \sqrt{x^2} = |x|.
\]
\textbf{Attention :} le résultat est \( |x| \), pas \( x \) (sauf si \( x \ge 0 \)).

\subsection*{\color{myred} 2. Exercice d’application }
Calculer \( (g \circ f)(x) \) pour :
\[
f(x) = \frac{1}{x} \quad (x \neq 0) \quad \text{et} \quad g(x) = 2x + 3.
\]
\textbf{Réponse :}
\[
(g \circ f)(x) = g\left(\frac{1}{x}\right) = \frac{2}{x} + 3.
\]

\newpage

\section*{\color{myred} II. Calcul de \( (f \circ g)(x) \) }

\subsection*{\color{myred} 1. Exemples }

\subsubsection*{\color{myred} a. Exemple 1 (mêmes fonctions) }
Reprenons 
\[
f(x) = 2x + 1 \quad \text{et} \quad g(x) = 3x - 4.
\]
Calculons \( (f \circ g)(x) \) :
\[
(f \circ g)(x) = f(g(x)) = f(3x-4) = 2(3x-4) + 1 = 6x - 8 + 1 = 6x - 7.
\]
\textbf{Comparaison :}
\[
(g \circ f)(x) = 6x - 1 \quad \text{et} \quad (f \circ g)(x) = 6x - 7.
\]
Les deux expressions sont différentes. On en déduit que la composition \textbf{n'est pas commutative}.

\subsubsection*{\color{myred} b. Exemple 2 (racine et carré) }
Soit 
\[
f(x) = x^2 \quad \text{et} \quad g(x) = \sqrt{x} \quad (x \ge 0).
\]
Alors :
\[
(f \circ g)(x) = f(g(x)) = f(\sqrt{x}) = (\sqrt{x})^2 = x \quad (x \ge 0).
\]

\subsection*{\color{myred} 2. Exercice d’application }
Calculer \( (f \circ g)(x) \) pour :
\[
f(x) = \frac{1}{x} \quad \text{et} \quad g(x) = 2x + 3 \quad \text{avec } x \neq -\frac{3}{2}.
\]
\textbf{Réponse :}
\[
(f \circ g)(x) = f(2x+3) = \frac{1}{2x+3}.
\]

\newpage

\subsubsection*{Synthèse et comparaison}
\begin{itemize}
    \item En général, \( g \circ f \neq f \circ g \).
    \item L'ordre de lecture est important : \( g \circ f \) se lit \og{} g rond f \fg{} et signifie \og on applique f puis g \fg.
\end{itemize}

\subsubsection*{Exercices de renforcement}
\begin{enumerate}
    \item Soit \( f(x) = x - 5 \) et \( g(x) = x^2 + 1 \). Calculer \( (g \circ f)(x) \) et \( (f \circ g)(x) \).
    \item Soit \( f(x) = \dfrac{x+1}{x-1} \) et \( g(x) = \dfrac{1}{x} \). Calculer les deux compositions (préciser les domaines de définition).
\item Choisir la bonne réponse et justifier le choix fait

\begin{table}[h]
\centering
\begin{tabular}{|m{7cm}|m{6cm}|}
\hline
 Si $f$ et $g$ sont les fonctions numériques définies respectivement par $f(x) = 2x^2 + 1$ et $g(x) = \dfrac{x-2}{x+1}$, alors, pour tout réel $x$, on a : 
& 
\begin{enumerate}
    \item[a.] $(g \circ f)(0) = 9$.
    \item[b.] $(g \circ f)(x) = 2\left(\dfrac{x-2}{x+1}\right)^2 + 1$.
    \item[c.] $(g \circ f)(x) = \dfrac{2x^2 - 1}{2(x^2 + 1)}$.
\end{enumerate} \\
\hline
\end{tabular}
\end{table}
    
\end{enumerate}

\section*{Bilan}
\begin{tcolorbox}[colback=gray!5!white, colframe=black!75!black]
\begin{itemize}
    \item La composition est une opération sur les fonctions.
    \item Elle n'est pas commutative.
    \item Pour calculer, on remplace la variable de la fonction extérieure par l'expression de la fonction intérieure.
\end{itemize}
\end{tcolorbox}

\section*{Évaluation formative (optionnelle)}
Donner deux fonctions au choix et demander de calculer les deux compositions à l'oral ou par écrit.

\end{document}